
retrouve la notion «~géométrique~» classique d'équivalence algébrique, envisagée dans 2.4.1. Aussi nous préférons dire qu'un élément de \(K^\bullet(T)\) est algébriquement équivalent à zéro si son image dans \(K^\bullet(T\otimes_k\bar k)\) l'est, où \(\bar k\) est une clôture algébrique de \(k\), et de procéder de même pour la définition de l'équivalence algébrique dans les groupes \(K_\bullet(T)\), \(\Gr^\bullet(T)\) et \(\Gr_\bullet(T)\).

\textbf{7.12.} Ces définitions étant posées, je dis que \emph{l'homomorphisme de spécialisation} (7.3.1) \emph{transforme éléments algébriquement équivalents à zéro en éléments algébriquement équivalents à zéro}, et définit par suite par passage au quotient un homomorphisme
\begin{equation}\tag{7.12.1}
  K_\bullet^{\mathrm{alg}}(X_t)\longrightarrow K_\bullet^{\mathrm{alg}}(X_s).
\end{equation}

Pour le voir, on se ramène immédiatement au cas où \(A\) est complet et \(k(s)\) algébriquement clos (moyennant une compatibilité facile explicitée dans 7.15). Il suffit de voir que \(\sigma_\bullet\) transforme un élément de la forme (7.11.1) de \(K_\bullet(X_t)\) en un élément algébriquement équivalent à zéro dans \(K_\bullet(X_s)\). Ici \(T=X_t\), et nous écrirons en conséquence \(C_t\), \(E_t\) au lieu de \(C\), \(E\) comme plus haut. On sait que \(C_t\) est projective sur \(k(t)\), donc (EGA IV 2.8.5) il existe un schéma projectif et plat \(C\) sur \(S\) dont la fibre générique soit \(C_t\). Utilisant la résolution des singularités des surfaces arithmétiques de Abhyankar [6], on voit qu'on peut même choisir pour \(C\) un schéma régulier. (N.B. La platitude est conservée par résolution, grâce à EGA IV 14.3.8). Comme \(C\) est propre sur \(S\), les points \(x_{0t}\) et \(x_{1t}\) de \(C_t\) rationnels sur \(k(t)\) se prolongent en des sections \(x_0\), \(x_1\) de \(C\) sur \(S\), et comme \(C\) est régulier, il est bien connu que \(C\) est lisse sur \(S\) en les points de \(x_0(S)\), \(x_1(S)\) (EGA IV, 17.12.1 c) \(\Rightarrow\) a) et 19.1.1). Par
